\documentclass[a4paper]{article}

% ---------- Encoding, lingua, font ----------
\usepackage[utf8]{inputenc}
\usepackage[T1]{fontenc}
\usepackage{textcomp}
\usepackage[italian]{babel}
\usepackage{xcolor}
\renewcommand{\familydefault}{\sfdefault}
\emergencystretch=1.5em

% ---------- Colori (UNA sola volta, prima di tikz/tcolorbox) ----------
\usepackage[dvipsnames,table]{xcolor}

% ---------- Matematica ----------
\usepackage{amsmath, amssymb, amsfonts, amsthm}
\usepackage{mathtools}
\DeclarePairedDelimiter\ceil{\lceil}{\rceil}
\DeclarePairedDelimiter\floor{\lfloor}{\rfloor}
\usepackage{cancel}
\usepackage{gensymb}

% ---------- Tabelle, figure, utilità ----------
\usepackage{booktabs, xltabular, tabularx, makecell}
\usepackage{subcaption}
\usepackage{float}
\usepackage{graphicx}
\usepackage{import}
\usepackage{pdfpages}
\usepackage{xifthen}
\pdfminorversion=7
\pdfsuppresswarningpagegroup=1
\graphicspath{{./figures/}}

\renewcommand{\arraystretch}{1.5}
\newcolumntype{Y}{>{\raggedright\arraybackslash}X}

% ---------- TikZ / pgfplots ----------
\usepackage{tikz}
\usetikzlibrary{shapes, arrows, automata, petri, decorations.markings,
                decorations.pathreplacing, positioning, calc, quotes}
\usepackage{circuitikz}
\usepackage{forest}
\usepackage{pgfplots}
\pgfplotsset{compat=1.18, width=10cm}

\tikzset{
    place/.style={circle, thick, draw=black, minimum size=6mm},
    state/.style={circle, thick, draw=black, fill=white, minimum size=6mm},
}

% ---------- Palette (solo quella nuova) ----------
\definecolor{thmcol}{HTML}{2563EB}   % blu     - teoremi
\definecolor{lemcol}{HTML}{7C3AED}   % viola   - lemmi
\definecolor{defcol}{HTML}{059669}   % verde   - definizioni
\definecolor{excol}{HTML}{EA580C}    % arancio - esempi
\definecolor{remcol}{HTML}{6B7280}   % grigio  - osservazioni

% ---------- Listings (riusa la palette nuova) ----------
\usepackage{listings}
\lstdefinestyle{mystyle}{
    backgroundcolor=\color{remcol!8},
    commentstyle=\color{defcol},
    keywordstyle=\color{thmcol},
    numberstyle=\tiny\color{remcol},
    stringstyle=\color{lemcol},
    basicstyle=\ttfamily\footnotesize,
    breakatwhitespace=false, breaklines=true, captionpos=b,
    keepspaces=true, numbers=left, numbersep=5pt,
    showspaces=false, showstringspaces=false, showtabs=false, tabsize=2
}
\lstdefinestyle{terminal}{
    backgroundcolor=\color{black},
    basicstyle=\ttfamily\color{green},
    frame=single, breaklines=true, showstringspaces=false
}
\lstset{style=mystyle}

% ---------- Teoremi, definizioni, esempi (tcolorbox) ----------
\usepackage[most]{tcolorbox}
\tcbuselibrary{theorems}

\tcbset{
  thmstyle/.style={
    enhanced, breakable, sharp corners,
    boxrule=0pt, frame hidden,
    borderline west={3pt}{0pt}{#1},
    colback=#1!6,
    colbacktitle=#1!14, coltitle=#1!60!black,
    fonttitle=\bfseries,
    titlerule=0pt,
    separator sign none,
    description delimiters parenthesis,
    before skip=12pt, after skip=12pt,
    left=8pt, right=8pt, top=4pt, bottom=4pt,
    overlay unbroken and first={\tikzset{external/export=false}},
    overlay middle and last={\tikzset{external/export=false}},
  }
}

\newtcbtheorem[number within=section]{theorem}{Teorema}{thmstyle=thmcol}{thm}
\newtcbtheorem[use counter from=theorem]{lemma}{Lemma}{thmstyle=lemcol}{lem}
\newtcbtheorem[use counter from=theorem]{definition}{Definizione}{thmstyle=defcol}{def}
\newtcbtheorem[use counter from=theorem]{example}{Esempio}{thmstyle=excol}{ex}
\newtcolorbox{remark}[1][]{thmstyle=remcol, title={Osservazione}, #1}

% ---------- Esercizi (mdframed) ----------
\usepackage{mdframed}
\theoremstyle{definition}
\newmdtheoremenv[
    linecolor=gray, leftmargin=0, rightmargin=0,
    innertopmargin=8pt, innerbottommargin=8pt
]{exercise}{Esercizio}[section]

% ---------- Figure svg ----------
\newcommand{\incfig}[1]{%
    \def\svgwidth{\columnwidth}
    \import{./figures/}{#1.pdf_tex}
}

% ---------- Hyperref (sempre per ultimo, prima di externalize) ----------
\usepackage{hyperref}
\hypersetup{colorlinks=false}

% ---------- Externalize ----------
\usetikzlibrary{external}
\tikzexternalize[prefix=figures/tikz/, optimize=false]
\tikzexternaldisable


\begin{document}
\begin{titlepage}
	\begin{center}
		\vspace*{1cm}

		\Huge
		\textbf{Automated Reasoning}

		\vspace{0.5cm}
		\LARGE
		UniVR - Dipartimento di Informatica

		\vspace{1.5cm}

		\textbf{Tommaso Toffali}

		\vfill


		\vspace{0.8cm}


		1° Semestre 2026/2027
    \end{center}
\end{titlepage}

\tableofcontents
\pagebreak
\part{}
\section{Introduction}
\subsection{Why study Automated Reasoning?}
Automated Reasoning is a foundamental field of Artificial Intelligence concerned with the design
and implementation of computionals procedures that enable machines to solve problems formulated in logical
languages. This procedures involve \textbf{logical inference and search}, allowing system to reason
in a deductive manner throught mechanical process.
\vspace{1em}
\noindent
\newline
Since the inception of Artificial intelligence at the Dartmouth Conference in 1956, one of the core challenges
has been to demonstrate machine intelligence throught the automated proving of mathematical and logical theorems.
This capability represents a cornerstone of symbolic AI, where knowledge is rapresented symbolically and 
stored in computer memory, and systematic procedures are applied to solve complex problems.

\begin{theorem}{Theorem Proving}
    Traditionally require human-level mathematical reasoning and insight. The goal of automated 
    reasoning is to mecanize this process, creating algorithms and systems capable of performing 
    logical deduction without human intervention.
\end{theorem}
\noindent
A logical language consist of formal symbols and syntactic rules. Within the symbolic approach to 
artificial intelligence, knowledge is encoded using these symbols, which are stored in computer memory.
The system then applies a sequence of mechanical steps to manipulate these symbols and arrive at solutions
to given problems.
\section{Historical Foundations and Theoretical Limits}
\subsection{The Hilbert-Turing Perspective}
The foundations of automated reasoning can be traced to fundamental questions in mathematical
logic and computability. Since Alan Turing's seminal work on Turing machines in 1936, researchers have
pursued David Hilbert's 1900 challenge: is it possible to develop a mechanical
procedure-executed deterministically by human computor-capable of solving any mathematical problem?
\vspace{1em}
\noindent 
\newline
\textbf{The Entscheidungsproblem}: A central question that emerged from this pursuit asks whether it is 
possible to define a mechanical method to detrmine the validity of formulas in first-order logic.
\vspace{1em}
\noindent 
\newline
Turing demostrated that the halting problem is undecidable using Turing machines. This fundamental result can be reduced to the validity
problem in first-order logic, providing a negative answer to Hilbert's challenge. This undecidability result forms one of the core theoretical
foundations of automated reasoning, estabilishing fundamental limits on what can be mechanically computed.

\subsection{Practical Applications of Automated Reasoning}
Automated reasoning find application across multiple domains requiring formal verification and systematic problem-solving:
\begin{itemize}
	\item \textbf{Planning}: Given an initial state, goal state, and formal model of agent actions, determine a sequence of moves that transforms the initial state into the goal state (utilizing SAT solvers)
	\item \textbf{Scheduling}: Generate temporal arrangements of tasks that stisfy complex temporal and resource constraints 
	\item \textbf{Software Analysis and Verfication}: Formal methods for ensuring program correctenss, including:
	\begin{itemize}
		\item Provable privacy guarantees in security protocols 
		\item Verification of distributed algorithms
		\item Analysis of randomized algorithms and their probabilistic properties
	\end{itemize}
	\item \textbf{Hybrid Approaches}: Integration of automated reasoning with machine learning techniques for enhanced problem-solving capabilities
\end{itemize}
These applications demonstrate the practical utility of automated reasoning in solving complex real-world problems that require rigouros logical 
analysis and systematic exploration of solution spaces.
\section{Core Problems in Automated Reasoning}
Automated reasoning fundamentally addresses three interconnected classes of decision problems in mathematical logic, each representing a distinct 
challenge in determining the truth properties of logical formulas.
\[
	\underline{\textbf{The Three Core Problems in Automated Reasoning}}
\]
\begin{enumerate}
	\item \textbf{Satisfiability (SAT)}: Does there exist a model?
	\item \textbf{Validity}: Is the formula true in all models?
	\item \textbf{Logical Consequence}: Does the conclusion follow from premises?
\end{enumerate}
\subsection{Satisfiability problems}
A satisfiability problem (SAT) concerns determining whether a given logical formula has a truth assignment that makes it true. Formally,
for a formula \( \phi \) with free variables, we seek to determine if there exists an interpretation \( \mathcal{I} \) such that \( \mathcal{I} \models \phi \). This 
represent the existence problem - does there exists a model making the formula true?
\subsection{Validity problems}
In contrast to satisfiability, validity problems ask whether a formula holds true under all possible interpretations.
A formula \( \phi \) is valid (denoted \( \models \phi \)) if every possible truth assignment to its variables result in 
\( \phi \) evaluating to true. This correspond to the unversal quantification over all possible models - the formula must 
hold in every conceivable scenario.
\subsection{Logical Consequence Problems}
The logical consequence problem represents a generalization of validity, asking whether a conclusion necessarly follows from 
given permises. Given a set of assumptions \( \varGamma = \left\{ \phi_1, \phi_2, \ldots, \phi_n \right\} \) and a conjecture 
\( \psi \), we determine if \( \varGamma \models \psi \) (i.e., if \( \psi \) is true in every model where all formulas 
in \( \varGamma \) are true). This captures the essence of logical deduction - what conclusions are forced by the given information.
\vspace{1em}
\noindent 
\newline
These three problems form the foundation of automated reasoning systems, with satisfiability 
serving as the most tractable and finding widespread application in pratical problem-solving contexts.
\section{Formal Language and Semantics}
\subsection{Logical Language: First-Order Logic}
The formal language employed in automated reasoning is first-order logic (FOL), a substantial extension of propositional logic
that incorporates quantification over individual elements.
\[
  \underline{\textit{First-Order Logic Components}}
\] 
\textit{First-order logic achieves greater expressive power by introducting}:
\begin{itemize}
	\item \textit{Individual variables} ranging over domain elements
	\item \textit{Function symbols} for constructing complex terms 
	\item \textit{Predicate symbols} for expressing relations between elements 
	\item \textit{Quantifiers} \( \forall, \exists \) enabling statements about all or some elements
\end{itemize}
This enhanced expressiveness allows first-order logic to formalize mathematical structures, relationship, and properties 
in a manner impossible in propositional logic, which can only express truth-functional combinations of atomic propositions 
\subsection{Signature and Universe of Discourse}
A signature \( \Sigma \) provied the vocabulary for interpreting formulas within a specific domain. Formally, a signature is 
a triple \( \Sigma = \langle \mathcal{S}, \mathcal{F}, \mathcal{P} \rangle \) where:
\begin{itemize}
	\item \( \mathcal{S} \) is a set of sorts (disjoint domains of interpretation)
	\item \( \mathcal{F} \) is a set of function symbols, each with an arity specifying the number of arguments 
	\item \( \mathcal{P} \) is a set of predicate symbols, which we use to write formulas, each with an associated arity 
\end{itemize}
The signature defines the syntactic elements available for constructing formulas, while the universe of discourse (or domain) provides 
the semantic interpretation. Sorts correspond to types in programming languages, enabling the formal specification of heterogeneous 
domains with different categories of objects.
\vspace{1em}
\noindent
\newline
The signature \( \Sigma = \langle \mathcal{S}, \mathcal{F}, \mathcal{P} \rangle \) cna be further decomposed to provide a complete syntactic 
framework. The set of function symbols \( \mathcal{F} \) is typically partitioned into two distinct subset:
\begin{itemize}
	\item \textbf{Constant symbols} \( \mathcal{C} \subseteq \mathcal{F} \): These are function symbols of arity 0, representing 
		named individuals in the domain. Constants have fixed interpretations and serve as proper names for specific domain elements 
	\item \textbf{Proper function symbols} \( \mathcal{F} \setminus \mathcal{C} \): These represent actual functions mapping tuples 
		of domain elements to domain elements, each with arity \( n \geq 1 \).
\end{itemize}
Formally, a function symbol \( f \in \mathcal{F} \) has type \( s_1 \times s_2 \times \ldots \times s_n \rightarrow s \), where 
\( s_1, \ldots, s_n, s \in \mathcal{S} \) are sorts indicating the domain and codomain of the function. For example:
\begin{itemize}
	\item \( c :\rightarrow \texttt{int} \) (a costant symbol denoting a specific integer)
	\item \( p :\rightarrow \texttt{color} \) (a costant symbol denoting a specific color)
	\item \( f : \texttt{int} \rightarrow \texttt{int} \) (a binary function like addition)
\end{itemize}
This decomposition enables precise specification of both individual elements (via constants) and functional relationships 
(via proper function symbol) within the formal language.
\vspace{1em}
\noindent
\newline
This complete signature structure \( \langle \mathcal{S}, \mathcal{C} \cup (\mathcal{F} \setminus \mathcal{C}) \cup \mathcal{P} \rangle \) provides 
the full syntactic foundation for expressing complex logical statements about structured domains.
\vspace{1em}
\noindent
\newline
Constant can be viewed as function symbols of arity - functions that require no arguments and directly denote specific domain elements.
Additionally, to construct quantified formulas and express generality, we require a countably infinte set of variables \( X = \left\{x,y,z,\ldots \right\} \) disjoint
from the signature symbols.
\vspace{1em}
\noindent
\newline
Thus, the function symbols may be comprehensively categorized as \( \mathcal{F} = \mathcal{C} \cup ( \mathcal{F} \setminus \mathcal{C}) \cup \mathcal{P} \), where 
\( \mathcal{P} \) represents the set of predicate symbols. Each predicate symbol \( p \in \mathcal{P} \) has arity \( n \geq 0 \) and 
type \( s_1 \times s_2 \times \ldots \times s_n \), where the final sort represents the truth value (typically a boolean sort).
\vspace{1em}
\noindent
\newline
Predicate symbols differ fundamentally from function symbols in that they represent relations 
rather than functions - they evaluate to truth values rather than domain elements. For example:
\begin{itemize}
	\item \textit{even}: \( \texttt{int} \rightarrow \texttt{bool} \) (unary predicate testing parity)
	\item \textit{less}: \( \texttt{int} \times \texttt{int} \rightarrow \texttt{bool} \) (binary relation for ordering)
\end{itemize}
The complete singature structure \( \langle \mathcal{S}, \mathcal{C} \cup (\mathcal{F} \setminus \mathcal{C}) \cup \mathcal{P} \rangle \) provides the full syntactic foundation for expressing 
complex logical statements about structured domains.
\vspace{1em}
\noindent
\newline
Constants can be viewed as function symbols of arity 0 - functions that require no arguments and directly denote 
specific domain elements. Additionally, to construct quantified formulas and express generality, we require a countably 
infinite set of variables \( X = \left\{ x,y,z, \ldots \right\} \) disjoint from the signature symbols.

\section{Syntactic Foundations: Terms and Formulas}
\subsubsection{Terms: The building blocks}
The most foundamental syntactic objects in first-order logic are \textbf{terms}, which represent individuals in the 
domain of discourse. The set of terms \( \mathcal{T}_{\Sigma} (\mathcal{X}) \) over signature \( \Sigma \) and variables \( \mathcal{X} \) is 
inductively defined as the smallest set satisying:
\begin{itemize}
	\item \textbf{Constants}: If \( c \in \mathcal{C} \) is a constant symbol, then \( c \in \mathcal{T}_{\Sigma} (\mathcal{X}) \).
	\item \textbf{Variables}: If \( x \in \mathcal{X} \) is a variable, then \( x \in \mathcal{T}_{\Sigma} (\mathcal{X}) \).
	\item \textbf{Function application}: If \( f \in \mathcal{F} \) is a function symbol of arity \( n \geq 1 \) and \( t_1, \ldots, t_n \in \mathcal{T}_{\Sigma} (\mathcal{X}) \) are terms, then \( f(t_1,\ldots,t_n) \in \mathcal{T}_{\Sigma} (\mathcal{X}) \).
\end{itemize}
Terms can be conceptualized as finite trees where: 
\begin{itemize}
	\item Leaf nodes are constants or variables 
	\item Internal nodes are function symbols applied to subterms
\end{itemize}
For example, the term \( f(x,g(a),z) \) has the tree representation:
\begin{center}
	\begin{tikzpicture}[node distance = 0.9cm]
		\node (f) at (0,1.6) {$f$};
		\node (x) at (-1,0.8) {$x$};
		\node (g) at (0,0.8) {$g$};
		\node (z) at (1,0.8) {$z$};
		\node (a) at (0,0) {$a$};

		\draw[blue] (f.south) -- (g.north);
		\draw[purple] (g.south) -- (a.north);
	\end{tikzpicture}
\end{center}
Terms may contain subterms, which correspond to subtrees in this representation. Complex terms are constructed 
recursively from simpler subterms, reflecting the hierarchical nature of the syntactic structure.
\subsubsection{From terms to Atomic Formulas}
Having established the notion of terms as representations of individuals, we now introduce atomic 
formulas to express relationships between these individuals. An \textbf{atomic formula} (or \textbf{atom}) is 
constructed by applying a predicate symbol to an appropriate number of terms.
\vspace{1em}
\noindent
\newline
Formally, if \( P \in \mathcal{P} \) is a predicate symbol of arity \( n \geq 0 \) and \( t_1, \ldots, t_n \in \mathcal{T}_{\Sigma} (\mathcal{X}) \) are terms,
then \( P(t_1,\ldots,t_n) \) is an atomic formula. When \( n = 0 \), the atomic formula reduces to the predicate 
symbol itself \( P \), which can be interpreted as a propositional constant.
\vspace{1em}
\noindent
\newline
The truth value of an atomic formula \( P(t_1,\ldots,t_n) \) depends on the interpretation of the predicate 
\( P \) and the denotations of the terms \( t_1,\ldots,t_n \) in a given model.
\subsection{Literals: Positive and Negative Atomic Formulas}
A \textbf{literal} is either an atomic formula (positive literal) or its negation (negative literal).
Formally, if \( \phi \) is an atomic formula, then:
\begin{itemize}
	\item \( \phi \) is a positive literal 
	\item \( \lnot \phi \) is a negative literal
\end{itemize}
Literals represent the basic building blocks for constructing complex logical formulas 
and play a crucial role in automated reasoning system, particularly in resolution-based theorem proving.
\subsection{Logical Connectives}
Logical connectives provide the means to construct complex formulas simpler ones. These operators
enable the expression of compound logical statements through systematic combination of atomic formulas and literals.
\vspace{1em}
\noindent
\newline
The primary logical connectives include:
\begin{itemize}
	\item \textbf{Negation}(\( \lnot \)): Express logical denial or contradiction 
	\item \textbf{Conjunction}(\( \wedge\)): Represents logical conjunction (AND) 
	\item \textbf{Disjunction}(\( \vee \)): Represent logical disjunction (OR)
	\item \textbf{Implication}(\( \rightarrow \)): Expresses logical implication (IF...THEN)
	\item \textbf{Biconditional}(\( \leftrightarrow \)): Represent logical equivalence (IF AND ONLY IF) 
\end{itemize}
\subsection{Formulas: The Complete Syntax}
The set of \textbf{formulas} (or \textbf{well-formed formulas}) \( \mathcal{F}_{\Sigma} (\mathcal{X}) \) over signature \( \Sigma \) and variables \( \mathcal{X} \) is 
inductively defined as the smallest set satisying:
\begin{itemize}
	\item \textbf{Atomic formulas}: If \( P \in \mathcal{P} \) is a predicate symbol of arity \( n \geq 0 \) and \( t_1, \ldots, t_n \in \mathcal{F}_{\Sigma} (\mathcal{X}) \), then \( P(t_1, \ldots, t_n) \in \mathcal{F}_{\Sigma} (\mathcal{X}) \).
	\item \textbf{Negation}: If \( \phi \in \mathcal{F}_{\Sigma} (\mathcal{X}) \), then \( (\lnot \phi) \in \mathcal{F}_{\Sigma} (\mathcal{X}) \). 
	\item \textbf{Binary connectives}: If \( \phi, \psi \in \mathcal{F}_{\Sigma} (\mathcal{X}) \), then:
	\begin{itemize}
		\item \( (\phi \wedge \psi) \in \mathcal{F}_{\Sigma} (\mathcal{X}) \) (conjunction)
		\item \( (\phi \vee \psi) \in \mathcal{F}_{\Sigma} (\mathcal{X}) \) (disjunction)
		\item \( (\phi \rightarrow \psi) \in \mathcal{F}_{\Sigma} (\mathcal{X}) \) (implication)
		\item \( (\phi \leftrightarrow \psi) \in \mathcal{F}_{\Sigma} (\mathcal{X}) \) (biconditional)
	\end{itemize}
\end{itemize}
This inductive definition ensures that all syntactically valid logical expressions can be systematically constructed 
from the basic building blocks of atomic formulas through the application 
of logical connectives.
\subsection{Quantifiers: Universality and Existence}
To complete the syntax of first-order logic, we introduce quantifiers that express generality 
and existence over the domain of discourse:
\begin{itemize}
	\item \textbf{Universal quantifier} \( (\forall) \): If \( \phi \in \mathcal{F}_{\Sigma} (\mathcal{X}) \) is a formula and \( x \in \mathcal{X} \) is a variable, then 
		\( (\forall x \phi) \in \mathcal{F}_{\Sigma} (\mathcal{X}) \).
	\item \textbf{Existential quantifier} \( (\exists) \): If \( \phi \in \mathcal{F}_{\Sigma} (\mathcal{X}) \) is a formula and \( x \in \mathcal{X} \) is a variable, then 
		\( (\exists x \phi) \in \mathcal{F}_{\Sigma} (\mathcal{X}) \)
\end{itemize}
Quantifiers bind variables within thei scope, enabling the expression of properties that hold 
for all elements of the domain (universal quantification) or for at least one element (existential quantification).
\subsection{Illustrative Example: Geometric Axioms}
Consider the following first-order formula expressing that any two distinct points determine a unique line:
\[
  \forall x \forall y ((P(x) \wedge P(y) \wedge x \neq y) \rightarrow \exists z (L(z) \wedge Q(z,x,y) \wedge \forall w (L(w) \wedge Q(w,x,y) \rightarrow w = z)))
\]
Where:
\begin{itemize}
	\item P(x) denotes that \( x \) is a point 
	\item L(z) denotes that \( z \) is a line 
	\item Q(z,x,y) denotes that line z passes through points and \( x \) and \( y \).
	\item \( x \neq y \) expresses that points \( x \) and \( y \) are distinct
\end{itemize}
The axioms formalizes the geometric principle that given any two distinct points in space, there 
exists exactly one line passing through both points, capturing a foundamental property of Euclidean geometry within the 
framework of first-order logic.
\subsection{Closed Formulas and Variable Scope}
A \textbf{closed formula} (or \textbf{sentence}) is a formula containing no free variables - all variables are bound by 
quantifiers. A variable is \textbf{free} ina a formula if it appears outside the scope of any 
quantifier that binds it; conversely, a variable is \textbf{bound} if it appears within the scope of a quantifier.
\vspace{1em}
\noindent
\newline
Formulas may contain both free and bound occurences of variables. For proper semantica interpretation, variables must be
\textbf{standardized apart}, meaning that no variable appears both free and bound 
in the same formula, and distinct quantifiers use distinct bound variables.
\part{Semantic Foundations}
\section{Interpretations and Model}
\subsection{The concept of Interpretation}
To assign meaning to syntactic formulas in first-order logic, we require a formal notion of 
\textbf{interpretation}. An interpretation \( \mathcal{I} \) provides a semantic mapping from the syntactic elements 
of a signature to concrete mathematical objects.
\begin{definition}(Interpretation Components)
\centering In the unsorted case, an interpretation consists of:
\begin{itemize}
	\item A non-empty set D called the \textbf{domain} (or universe of discourse)
	\item An interpretation function \( \Phi \) that maps:
	\begin{itemize}
		\item Each predicate symbol \( P \in \mathcal{P} \) of arity \( n \) to a subset \( \Phi(P) \subseteq D^n \) (representing an \( n \)-ary realtion over the domain)
		\item Each function symbol \( f \in \mathcal{F} \) of arity \( n \) to a function \( \Phi(f) : D^n \rightarrow D \).
		\item Each constants symbol \( c \in \mathcal{C} \) to a specific element \( \Phi(c) \in D \).
 	\end{itemize}
\end{itemize}
\end{definition}
The interpretation function \( \Phi \) bridges the gap between syntax and semantics, transforming abstract symbols into concrete 
mathematical entities.
\subsection{Example: Interpreting a Formula over Natural Numbers}
Consider the first-order formula:
\[
  \forall x (\lnot P(x) \rightarrow P(f(x)))
\]
The formula contains:
\begin{itemize}
	\item A unary predicate symbol \( P \) 
	\item A unary function symbol \( f \)
	\item A universally quantified variable \( x \)
\end{itemize}
Since the formula is closed (contains no free variables), its truth value depends entirely on the chosen interpretation.
\vspace{1em}
\noindent
\newline
\textbf{Constructing an interpretation}: Let us define an interpretation \( \mathcal{I} = \langle D, \Phi \rangle \) where:
\begin{itemize}
	\item \textbf{Domain}: \( D = \mathbb{N} \) (the set of natural numbers)
	\item \textbf{Function interpretation}: \( \Phi(f) : \mathbb{N} \rightarrow \mathbb{N} \) is the successor function, defined by \( \Phi(f)(n) = n + 1 \) for all \( n \in \mathbb{N} \).
	\item \textbf{Predicate interpretation}: \( \Phi(P) = \left\{0,2,4,6,\ldots \right\} = \left\{ 2k \;|\; k \in \mathbb{N} \right\} \) (the set of even natural numbers)
\end{itemize}
\subsection{Semantic Evaluation}
Under this interpretation \( \mathcal{I} \), the formula \( \forall x (\lnot P(x) \rightarrow P(f(x))) \) receives the following 
semantic reading:
\vspace{1em}
\noindent
\newline
\textit{For all natural numbers x, if x is not even, then the successor of x is even}.
\vspace{1em}
\noindent
\newline
To determine whether this formula is true in interpretation \( \mathcal{I} \) we must verify whether the implication holds for every element of the domain:
\begin{itemize}
	\item If \( x \) is odd (i.e., \( x \notin \Phi(P)) \), then \( x + 1 \) must be even (i.e., \( x + 1 \in \Phi(P)) \)
\end{itemize}
The statement is indeed true in the natural numbers: the successor of any odd number is even.
Therefore, the formula is \textbf{satisfied} by interpretation \( \mathcal{I} \), and we write \( \mathcal{I} \models \forall x(\lnot P(x) \rightarrow P(f(x))) \).
\subsection{Formal Verification of Satisfaction}
The formally verify that \( \mathcal{I} \models \forall x (\lnot P(x) \rightarrow P(f(x))) \), we must check that for all \( n \in \mathbb{N} \), the 
interpretation satisfies the implication with the substitution \( \left\{ x \leftarrow n \right\} \).
\vspace{1em}
\noindent
\newline
By the semantics of implications, \( \mathcal{I} \) satisfies \( (\lnot P(x) \rightarrow P(f(x))) \) with assignment \( \left\{ x \leftarrow n \right\} \) if and only if:
\begin{itemize}
	\item Either \( \mathcal{I} \) does not satisfy \( \lnot P(x) \) with \( \left\{ x \leftarrow n \right\} \) (i.e., \( \mathcal{I} \) satisfies \( P(x) \) with \( \left\{ x \leftarrow n \right\} \)), or 
	\item \( \mathcal{I} \) satisfies \( P(f(x)) \) with \( \left\{ x \leftarrow n \right\} \)
\end{itemize}
Equivalently, for all \( n \in \mathbb{N} \), either \( n \in \Phi(P) \) or \( \Phi(f)(n) \in \Phi(P) \).
\vspace{1em}
\noindent
\newline 
\textbf{Truth condition for atomic formulas}: An atomic formula \( P(t) \) is true under interpretation 
\( \mathcal{I} \) if and only if the interpretation of the term \( t \) belongs to the interpretation of the predicate \( \Phi(P) \)..
\vspace{1em}
\noindent
\newline 
In our example, for all \( n \in \mathbb{N} \):
\begin{itemize}
	\item Either \( n \in \Phi(P) \) (meaning \( n \) is even), or 
	\item \( \Phi(f)(n) = n + 1 \in \Phi(P) \) (meaning the successor of \( n \) is even)
\end{itemize}
This condition holds for all natural numbers, confirming that the formula is satisfied by the interpretation.
\subsection{Formal Definition of Interpretation}
Having illustrated the concept through examples, we now provide the complete formal definition 
of an interpretation in first-order logic.
\vspace{1em}
\noindent
\newline
An \textbf{interpretation} \( \mathcal{I} \) consists of a triple \( \mathcal{I} = \langle D, \Phi, \beta \rangle \) where:
\begin{itemize}
	\item \textbf{Domain}: \( D \) is a non-empty set representing the universe of discourse 
	\item \textbf{Interpretation function} \( \Phi \) maps signature symbols to semantic objects:
	\begin{itemize}
		\item For each constant symbol \( c \in \mathcal{C}: \Phi(c) = d \in D \).
		\item For each function symbol \( f \in \mathcal{F} \) of arity \( n \geq 1: \Phi(f) : D^n \rightarrow D \), where \( D^n \) denotes the \( n \)-fold Cartesian product \( D \times D \times \ldots \times D \).
		\item For each predicate symbol \( P \in \mathcal{P} \) of arity \( n \geq 0: \Phi(P) \subseteq D^n \) (an \( n \)-ary relation over \( D \)) 		
	\end{itemize}
	\item \textbf{Variable assignment} \( \beta \) maps each variable to a domain element:
	\begin{itemize}
		\item \( \beta : \mathcal{X} \rightarrow D \) such that \( \beta(x) = d \in D \) for all variables \( x \in \mathcal{X} \)
	\end{itemize}
\end{itemize}
\textbf{Assignement update notation}: Given an assignment \( \beta \) and a substitution \( x \leftarrow d \), we define 
the updated assignment \( \beta[x \leftarrow d] \) as:
\[
  \beta[x \leftarrow d](y) = \begin{cases}
	d & \text{if } y = x \\
	\beta(y) & \text{otherwise}
  \end{cases}
\]
This notation allows us to formally express the evaluation of formulas with different variable bindings, 
which is essential for defining the semantics of quantified formulas.
\section{Satisfaction Relation}
\subsection{The satisfaction relation}
Having defined interpretations and variable assignments, we now formalize the central semantica concept: when 
does an interpretation \textbf{satisfy} a formula? The satisfaction relation, denoted 
\( \mathcal{I} \models_{\beta} \Phi \), specifies the conditions under which a formula \( \Phi \) is true in interpretation \( \mathcal{I} \) with 
variable assignment \( \beta \).
\vspace{1em}
\noindent
\newline
Let \( F \) and \( G \) denote arbitraty formulas. The satisfaction realtion is defined inductively on the 
structure or formulas.
\subsection{Interpretation of Terms}
Before defining satisfaction for formulas, we must specify how terms are interpreted. Given an interpretation \( \mathcal{I} = \langle D, \Phi, \beta \rangle \),
the interpretation of a term \( t \), denoted \( \Phi_{\beta}(t) \), is defined recursively:
\begin{itemize}
	\item \textbf{Constants}: \( \Phi_{\beta}(c) = \Phi(c) \in D \) 
	\item \textbf{Variables}: \( \Phi_{\beta}(x) = \beta(x) \in D \)   
	\item \textbf{Function application}: \( \Phi_{\beta}(f(t_1, \ldots, t_n)) = \Phi(f)(\Phi(t_1), \ldots, \Phi_{\beta}(t_n)) \in D \) 
\end{itemize}
Each term evalueates to a specific element of the domain under a given interpretation.
\subsection{Satisfaction for Atomic Formulas}
\textbf{Base case} (atomic formulas): An interpretation \( \mathcal{I} \) satisfies an atomic formula \( P(t_1,\ldots, t_n) \) under 
assignment \( \beta \) if and only if the tuple of interpreted terms belongs to the interpretation of the predicate:
\[
  \mathcal{I} \models_{\beta} P(t_1,\ldots,t_n) \text{ iff } (\Phi_{\beta}(t_1), \ldots, \Phi_{\beta}(t_n)) \in \Phi(P)
\]
\subsection{Satisfaction for Logical Connectives}
The satisfaction relation extends to compound fomrulas through the following inductive clauses. Note that all satisfaction relations 
are parametrized by the variable assignment \( \beta \), written as \( \models_{\beta} \)
\begin{itemize}
	\item \textbf{Negation}:
	\[
	  \mathcal{I} \models_{\beta} \lnot F \text{ iff } \mathcal{I} \nvDash_{\beta} F 
	\]
	\item \textbf{Conjunction}: 
	\[
	  \mathcal{I} \models_{\beta} (F \wedge G) \text{ iff } \mathcal{I} \models_{\beta} F \text{ and } \mathcal{I} \models_{\beta} G
	\]
	\item \textbf{Disjunction}:
	\[
	  \mathcal{I} \models_{\beta} (F \vee G) \text{ iff } \mathcal{I} \models_{\beta} F \text{ or } \mathcal{I} \models_{\beta} G
	\]
	\item \textbf{Implication}:
	\[
	  \mathcal{I} \models_{\beta} (F \rightarrow G) \text{ iff } \mathcal{I} \nvDash_{\beta} F \text{ or } \mathcal{I} \models_{\beta} G 
	\]
	Equivalently: if \( \mathcal{I} \models_{\beta} F \) then \( \mathcal{I} \models_{\beta} G \).
	\item \textbf{Biconditional}:
	\[
	  \mathcal{I} \models_{\beta} (F \leftrightarrow G) \text{ iff } \mathcal{I} \models_{\beta} F \text{ if and only if } \mathcal{I} \models_{\beta} G
	\]
	The formula \( \exists x G \) is satisfied if \( G \) is satisfied under at least one assignment to \( x \).
\end{itemize}
This completes the inductive definition of the satisfaction relation, providing a complete semantics for first-order logic.
\subsection{Example: Quantifier Evaluation}
Consider the formula \( \forall x \exists y R(x,y) \), which can be read as "for all \( x \), there exists \( y \) such that 
\( x < y \)".
\vspace{1em}
\noindent
\newline
Let \( \mathcal{I} = \langle \mathbb{N}, \Phi \rangle \) where \( \Phi(R) = \left\{ (n,m) \in \mathbb{N} \times \mathbb{N} \;|\; n < m \right\} \) (the 
standard ordering on natural numbers).
\vspace{1em}
\noindent
\newline
To verify \( \mathcal{I} \models \forall x \exists y R(x,y) \), we apply the semantic definitions:
\begin{align*}
	& \mathcal{I} \models \forall x\, \exists y\, R(x,y) \\
	\iff\ & \text{per ogni } n \in \mathbb{N},\
		\mathcal{I} \models_{\beta[x \leftarrow n]} \exists y\, R(x,y) \\
	\iff\ & \text{per ogni } n \in \mathbb{N},\ \text{esiste } m \in \mathbb{N} \text{ tale che }
		\mathcal{I} \models_{\beta[x \leftarrow n,\, y \leftarrow m]} R(x,y) \\
	\iff\ & \text{per ogni } n \in \mathbb{N},\ \text{esiste } m \in \mathbb{N} \text{ tale che }
		(n,m) \in \Phi(R) \\
	\iff\ & \text{per ogni } n \in \mathbb{N},\ \text{esiste } m \in \mathbb{N} \text{ tale che }
		n < m
\end{align*}
This statement is \textbf{false} in \( \mathbb{N} \) because there is no natural number greater than all natural numbers.
However, it would be true \( \mathbb{Z} \) or \( \mathbb{Q} \), demonstrating how truth depends on the chosen interpretaion.
\section{Semantic Properties of Formulas}
Having established the satisfaction relation, we now define fundamental semantic properties that classify formulas according 
to their to their truth behavior across different interpretations.
\begin{definition}{Semantic Properties of Formulas}
\centering
\textbf{Satisfiability}
\newline
A formula \( F \) is \textbf{satisfiable} if there exists an interpretation \( \mathcal{I} \) such that \( \mathcal{I} \models F \).
\[
  F \text{ is satisfiable } \Leftrightarrow \exists \mathcal{I} : \mathcal{I} \models F
\]	
\textbf{Validity}
\newline
A formula \( F \) is \textbf{valid} (or a \textbf{tautology}) if for all interpretations \( \mathcal{I} \), we have \( \mathcal{I} \models F \).
\[
  F \text{ is valid } \Leftrightarrow \forall \mathcal{I} : \mathcal{I} \models F 
\]
\textbf{Unsatisfiability}
\newline
A formula \( F \) is \textbf{unsatisfiable} if there exists no interpretation \( \mathcal{I} \) such that \( \mathcal{I} \models F \).
\[
  F \text{ is unsatisfiable } \Leftrightarrow \nexists \mathcal{I} : \mathcal{I} \models F  
\]
\textbf{Invalidaity}
\newline
A formula \( F \) is \textbf{invalid} if there exists an interpretation \( \mathcal{I} \) such that \( \mathcal{I} \nvDash F \).
\[
  F \text{ is invalid } \Leftrightarrow \exists \mathcal{I} : \mathcal{I} \nvDash F
\]
\end{definition}
\subsection{Satisfiability}
A formula \( F \) is \textbf{satisfiable} if there exists an interpretation \( \mathcal{I} \) such that \( \mathcal{I} \models F \). In thi case, 
we say that \( \mathcal{I} \) is a \textbf{model} of \( F \).
\subsection{Validity}
We denote validity by \( \models F \).
\subsection{Unsatisfiability}
Unsatisfiable formulas are false in every interpretation.
\subsection{Invalidity}
In this case, we say that \( \mathcal{I} \) is a \textbf{countermodel} (or \textbf{counterexample}) for \( F \).
\subsection{Relationship Between Semantic Properties}
These semantic properties are intimately related through negation:
\begin{theorem}{Validity and Unsatisfiability}
\centering
\textit{A formula F is valid if and only if \( \lnot F \) is unsatisfiable}.
\vspace{1em}
\noindent
\newline
\textit{Proof. Suppose  F is valid. Then for all interpretations \( \mathcal{I} \), we have \( \mathcal{I} \models F \), which implies \( \mathcal{I} \nvDash \lnot F \).}
\newline
Therefore, \( \lnot F \) is unsatisfiable.
\vspace{1em}
\noindent
\newline
Conversely, if \( \lnot F \) is unsatisfiable, then for all interpretations \( \mathcal{I} \), we have \( \mathcal{I} \nvDash \lnot F \), which implies 
\( \mathcal{I} \models F \). Therefore, \( F \) is valid.
\end{theorem}
\begin{theorem}{Satisfiability and Invalidity}
\centering
\textit{A formula F is satisfiable if and only if \( \lnot F \) is invalid}
\vspace{1em}
\noindent
\newline
\textit{Proof}. \( F \) is satisfiable if and only if there exists an interpretation \( \mathcal{I} \) such that \( \mathcal{I} \models F \), which 
is equivalent to \( \mathcal{I} \nvDash \lnot F \). This is precisely the condition for \( \lnot F \) to be invalid.
\end{theorem}
These relationships demonstrate the duality between validity and unsatisfiability, and between 
satisfiability and invalidity, providing a complete characterization of the semantic landscape of first-order logic.
\section{Logical Consequence}
\subsection{The Three Fundamental Query Types}
Automated reasoning system must be capable of resolving three fundamental types of semantic queries:
\begin{itemize}
	\item \textbf{Satisfiability queries}: Is formula \( F \) satisfiable?
	\item \textbf{Validity queries}: Is formula \( F \) valid?
	\item \textbf{Logical consequence queries}: Does \( F \) follow logically from a set of hypotheses \( H \)?
\end{itemize}
\subsection{Definition of Logical Consequence}
\begin{definition}{Logical Consequence}
\centering
Let \( H = \left\{ \phi_1, \phi_2, \ldots, \phi_n \right\} \) be a set of formulas (the \textbf{hypotheses} or \textbf{premises}) and let 
\( F \) be a formula (the \textbf{conclusion}). We say that \( F \) is a \textbf{logical consequence} of \( H \), written \( H \models F \), if:
\[
  H \models F \Leftrightarrow \text{ for all interpretations } \mathcal{I}, \text{ if } \mathcal{I} \models H \text{ then } \mathcal{I} \models F
\]
\end{definition}
In other words, \( F \) is a logical consequence of \( H \) if every model of \( H \) is also a model of \( F \). Here, \( \mathcal{I} \models H \)
means that \( \mathcal{I} \) satisfies all formulas in \( H \).
\subsection{The deduction Theorem}
The deduction theorem establishes a fundamental connection between logical consequence and validity, reducing the probem of checking entailment 
to the problem of checking validity.
\begin{definition}{Deduction theorem}
\centering 
\textit{Let H be a set of formulas and F be a formula}.
\newline
\textit{Then}:
\[
  H \models F \text{ iff } \models (H \rightarrow F)
\]
\textit{where \( H \rightarrow F \) denotes the implication from the conjunction of all formulas in H to F}
\end{definition}
\textit{Proof}. We prove both directions:
\newline
(\( \Rightarrow \)) Assume \( H \models F \). We must show that \( H \rightarrow F \) is valid, i.e., for all interpretations \( \mathcal{I} \), we have 
\( \mathcal{I} \models H \rightarrow F \).
\vspace{1em}
\noindent
\newline
Let \( \mathcal{I} \) be an arbitraty interpretation. By the semantics of implication, \( \mathcal{I} \models H \rightarrow F \), holds if 
either:
\begin{itemize}
	\item \( \mathcal{I} \nvDash H \) (the antecedent is false), or 
	\item \( \mathcal{I} \models F \) (the consequent is true)
\end{itemize}
If \( \mathcal{I} \nvDash H \), then \( \mathcal{I} \models H \rightarrow F \) trivially. If \( \mathcal{I} \models H \), then by the hypothesis
\( H \models F \), we have \( \mathcal{I} \models F \), so \( \mathcal{I} \models H \rightarrow F \). In both cases, \( \mathcal{I} \models H \rightarrow F \).
\vspace{1em}
\noindent
\newline
(\( \Leftarrow \)) Conversely, suppose \( H \rightarrow F \) is valid. Then for all interpretations \( \mathcal{I} \), we have \( \mathcal{I} \models H \rightarrow F \).
By the semantics of implications, for all \( \mathcal{I} \), either \( \mathcal{I} \nvDash H \) or \( \mathcal{I} \models F \). Therefore, for all 
\( \mathcal{I} \) such that \( \mathcal{I} \models H \), we must have \( \mathcal{I} \models F \), which the definition of \( H \models F \).
\subsection{Notational Ambiguity of \( \models \)}
The symbol \( \models \) is used in three distinct but related contexts:
\begin{itemize}
	\item \textbf{Satisfaction}: \( \mathcal{I} \models F \) means interpretation \( \mathcal{I} \) satisfies formula \( F \).
	\item \textbf{Validity}: \( \models F \) means formula \( F \) is valid (satisfied by all interpretations)
	\item \textbf{Entailment}: \( H \models F \) means formula \( F \) is a logical consequence of the set of formulas \( H \).
\end{itemize}
While this overloading may initially seem confusing, these three uses are smentically coherent and represent different aspects of the same fundamental 
semantic relation.
\subsection{Reduction to Unsatisfiability}
A fundamental result connects logical consequence to unsatisfiability, providing another characterization of entailment.
\begin{theorem}{Jey Theorem: Reduction to Unsatisfiability}
\centering
\textit{Let H be a set of formulas and F be a formula. Then:}
\[
  H \models F \text{ iff } H \cup \left\{ \lnot F \right\} \text{ is unsatisfiable}
\]
\end{theorem}
\textit{Proof}. We prove both directions:
\vspace{1em}
\noindent
\newline
\( (\Rightarrow) \) Assume \( H \models F \). We must show that \( H \cup \left\{ \lnot F \right\} \) is unsatisfiable.
\vspace{1em}
\noindent
\newline
Suppose for contradiction that there exists an interpretation \( \mathcal{I} \) such that \( \mathcal{I} \models H \cup \left\{ \lnot F \right\} \). Then 
\( \mathcal{I} \models H \), and \( \mathcal{I} \models \lnot F \), which implies \( \mathcal{I} \nvDash F \). However, sincer \( H \models F \) and \( \mathcal{I} \models H \),
we must have \( \mathcal{I} \models F \), which is a contradiction. Therefore, \( H \cup \left\{ \lnot F \right\} \) is unsatisfiable.
\vspace{1em}
\noindent
\newline
\( (\Leftarrow) \) Assume \( H \cup \left\{ \lnot F \right\} \) is unsatisfiable. We must show that \( H \models F \).
\vspace{1em}
\noindent
\newline
Let \( \mathcal{I} \) be any interpretation such that \( \mathcal{I} \models H \). Since \( H \cup \left\{ \lnot F \right\} \) is unsatisfiable,
there is no interpretation that satisfies both \( H \) and \( \lnot F \). Therefore, \( \mathcal{I} \nvDash \lnot F \), which implies \( \mathcal{I} \models F \). Since
this holds for all interpretations satisfying \( H \), we conclude that \( H \models F \).
\section{Refutational Reasoning}
\subsection{Proof by Contradiction}
\begin{definition}{Refutational reasoning (Proof by contradiction)}
\centering
\textbf{Refutational reasoning} (or \textbf{proof by refutation}) is a fundamental techinique in automated reasoning on the principle 
of \textit{reductio ad absurdum}. To prove that a conclusion \( F \) follows from hupotheses \( H \), we assume the negation of the conclusion 
and derive a contradiction.
\begin{enumerate}
	\item \textbf{Goal}: Prove \( H \models F \) (where \( H \) represents assumpions and \( F \) represents the conjecture) 
	\item \textbf{Method}: Show that \( H \cup \left\{ \lnot F \right\} \) is unsatisfiable 
	\item \textbf{Conclusion}: If \( H \cup \left\{ \lnot F \right\} \) is unsatisfiable, then \( H \models F \) by the previous theorem
\end{enumerate}
\end{definition}
\subsection{The refutational Framework}
The automated reasoning process can be formalized as follows:
\begin{center}
	\begin{tikzpicture}[>=stealth]
		% nodi principali
		\node (in) at (0,0) {$H \cup \{\neg F\}$};
		\node[draw] (dp) at (4,0) {Decision Procedure};
		\node (sat)   at (8,2.5)  {\textbf{Satisfiable}};
		\node (unsat) at (8,-2.5) {\textbf{Unsatisfiable}};
  
		% testi a destra
		\node[right=0.3cm of sat, text width=4.5cm, align=left]
	  	{Model found: counterexample to $H \models F$};
		\node[right=0.3cm of unsat, text width=4.5cm, align=left]
	  	{Proof found (refutation): $H \models F$};
  
		% frecce
		\draw[->] (in) -- (dp);
		\draw[->] (dp.north east) -- node[above left] {No}  (sat.south west);
		\draw[->] (dp.south east) -- node[below left] {Yes} (unsat.north west);
	\end{tikzpicture}
\end{center}
\textbf{Two possible outcomes}:
\begin{itemize}
	\item \textbf{Unsatisfiable}: The procedure derives a contradiction, producing a \textbf{refutation} (proof) that \( H \models F \).
	\item \textbf{Satisfiable}: The procedure finds a model \( \mathcal{I} \) of \( H \cup \left\{ \lnot F \right\} \), which serves as a \textbf{counterexample} showing that \( H \nvDash F \).
\end{itemize}
This refutational approach forms the foundation of modern automated theorem provers and SAT solvers, reducing the problem of proving logical consequence 
to the problem of detecting unsatisfiability.
\part{Free Variables and Satisfiability}
\section{Satisfiability of Formulas with Free Variables}
\subsection{The problem of Free Variables}
DA RECUPERARE FINO A SLIDE 39!!


\end{document}